% Generated by scripts/assemble.py; corrections belong in source/corrections.json.
\examhead{130}{武汉大学 2024 年数学分析试卷}{123}
% p124-b22; source PDF page 124
\par\Needspace{5\baselineskip}
\blockanchor{p124-b22}
\majorquestion{一}求导数或极限.\par
% p124-b23; source PDF page 124
\blockanchor{p124-b23}
\question{1}$\lim_{x\to 0}\frac{\sin x - \arctan x}{\tan x - \arcsin x}.$\par
% p125-b00; source PDF page 125
\blockanchor{p125-b00}
\question{2}$y = \frac{\arcsin x}{\sqrt{1 - x^2}}$ 求 $y^{(2023)}(0)$\par
% p125-b01; source PDF page 125
\blockanchor{p125-b01}
\question{3}$\lim_{n\to \infty}\sum_{k = 1}^{n}\sin \frac{k}{n}\cdot \sin \frac{k}{n^2}.$\par
% p125-b02; source PDF page 125
\par\Needspace{5\baselineskip}
\blockanchor{p125-b02}
\majorquestion{二}求积分.\par
% p125-b03; source PDF page 125
\blockanchor{p125-b03}
\question{1}以 $x = ty$ 参数化曲线 $x^{2} + y^{3} = xy$, 求曲线所围区域的面积.\par
% p125-b04; source PDF page 125
\par\Needspace{6\baselineskip}
\blockanchor{p125-b04}
\question{2}求曲面积分\par
\nopagebreak[4]
% p125-b05; source PDF page 125
\blockanchor{p125-b05}
\displayformula{\iint_ {S} (x y + y z + z x) \mathrm{d} S.}
% p125-b06; source PDF page 125
\blockanchor{p125-b06}
其中 $S$ 是 $z = \sqrt{x^2 + y^2}$ 被 $x^{2} + y^{2} = 2ay$ 截掉的部分.\par
% p125-b07; source PDF page 125
\blockanchor{p125-b07}
\question{3}用含参量积分求 $\int_0^1\frac{\arctan x}{x\sqrt{1 - x^2}}\mathrm{d}x.$\par
% p125-b08; source PDF page 125
\par\Needspace{5\baselineskip}
\blockanchor{p125-b08}
\majorquestion{三}证明题.\par
% p125-b09; source PDF page 125
\par\Needspace{6\baselineskip}
\blockanchor{p125-b09}
\question{1}证明：\par
\nopagebreak[4]
% p125-b10; source PDF page 125
\blockanchor{p125-b10}
\displayformula{\frac {2 n}{2 n + 1} \leq \frac {\left(1 + \frac {1}{n}\right) ^ {n}}{e} \leq \frac {2 n + 1}{2 n + 2} (n \geq 1).}
% p125-b11; source PDF page 125
\par\Needspace{6\baselineskip}
\blockanchor{p125-b11}
\question{2}设 $f(x)$ 在 $(0, +\infty)$ 上三次可导，且 $\lim_{x \to +\infty} f(x)$ 与 $\lim_{x \to +\infty} f'''(x)$ 均存在，证明：\par
\nopagebreak[4]
% p125-b12; source PDF page 125
\blockanchor{p125-b12}
\displayformula{\lim _ {x \rightarrow + \infty} f ^ {\prime} (x) = \lim _ {x \rightarrow + \infty} f ^ {\prime \prime} (x) = \lim _ {x \rightarrow + \infty} f ^ {\prime \prime \prime} (x) = 0.}
% p125-b13; source PDF page 125
\par\Needspace{6\baselineskip}
\blockanchor{p125-b13}
\question{3}设 $f(x)$ 是 $[-1, 1]$ 上的连续函数, 证明:\par
\nopagebreak[4]
% p125-b14; source PDF page 125
\blockanchor{p125-b14}
\displayformula{\lim _ {\varepsilon \rightarrow 0 ^ {+}} \int_ {- 1} ^ {1} \frac {\varepsilon f (x)}{\varepsilon^ {2} + x ^ {2}} \mathrm{d} x = \pi f (0).}
% p125-b15; source PDF page 125
\par\Needspace{6\baselineskip}
\blockanchor{p125-b15}
\question{4}设 $x > 0, y > 0$, 证明:\par
\nopagebreak[4]
% p125-b16; source PDF page 125
\blockanchor{p125-b16}
\displayformula{x y - e ^ {x - 1} \leq y \ln y.}
% p125-b17; source PDF page 125
\par\Needspace{6\baselineskip}
\blockanchor{p125-b17}
\question{5}证明: $\int_{0}^{1}(x\ln x)^{n}\mathrm{d}x = \frac{(-1)^{n}n!}{(n + 1)^{n + 1}}$，并利用此结论证明：\par
\nopagebreak[4]
% p125-b18; source PDF page 125
\blockanchor{p125-b18}
\displayformula{\int_ {0} ^ {1} x ^ {- x} \mathrm{d} x = \sum_ {n = 1} ^ {\infty} \frac {1}{n ^ {n}}.}
% p125-b19; source PDF page 125
\par\Needspace{6\baselineskip}
\blockanchor{p125-b19}
\question{6}设 $\{f_n(x)\}$ 是区间 $[a, b]$ 上一致收敛于 $f$ 的可积函数列. 证明: $f$ 在 $[a, b]$ 上可积, 且\par
\nopagebreak[4]
% p125-b20; source PDF page 125
\blockanchor{p125-b20}
\displayformula{\int_ {a} ^ {b} f (x) \mathrm{d} x = \lim _ {n \rightarrow \infty} \int_ {a} ^ {b} f _ {n} (x) \mathrm{d} x.}
